\documentclass[10pt,letterpaper]{article}
\usepackage[T1]{fontenc}
\usepackage[utf8]{inputenc}
\usepackage[spanish,es-tabla]{babel}
\usepackage{csquotes}
\usepackage{mathpazo}
\usepackage[scaled=0.9]{helvet}
\usepackage{courier}
\usepackage{calc}
\usepackage[top=2cm, bottom=2cm, left=2cm, right=2cm]{geometry}
\usepackage{xcolor}
\usepackage{booktabs}
\usepackage{longtable}
\usepackage{array}
\usepackage{ragged2e}
\usepackage{xurl}
\usepackage[strings]{underscore}
\usepackage{fancyhdr}
\usepackage{sectsty}
\usepackage{tocloft}
\usepackage[colorlinks=true,linkcolor=blue,urlcolor=blue,citecolor=blue,breaklinks=true]{hyperref}
\usepackage{enumitem}

\definecolor{uncpblue}{HTML}{1A237E}
\allsectionsfont{\color{uncpblue}}

\renewcommand{\cftsecleader}{\cftdotfill{\cftsecdotsep}}
\cftpagenumbersoff{section}
\cftpagenumbersoff{subsection}
\cftpagenumbersoff{subsubsection}
\renewcommand{\cftsecfont}{\normalfont}
\renewcommand{\cftsubsecfont}{\normalfont}
\renewcommand{\cftsubsubsecfont}{\normalfont}
\setlength{\parindent}{0pt}
\setlength{\parskip}{4pt}
\setlength{\tabcolsep}{3pt}
\renewcommand{\arraystretch}{0.85}

\pagestyle{fancy}
\fancyhf{}
\fancyhead[L]{\small\textcolor{uncpblue}{\textbf{SGSI-UNCP}}}
\fancyhead[R]{\small\textcolor{uncpblue}{R-SGSI-02}}
\fancyfoot[C]{\thepage}
\renewcommand{\headrulewidth}{0.4pt}
\renewcommand{\footrulewidth}{0.4pt}
\setlength{\headheight}{14pt}

\setlist{nosep,leftmargin=*,after=\vspace{2pt}}
\setlength{\emergencystretch}{2em}

\newcommand{\makecover}{
\begin{titlepage}
\centering
\vspace*{2cm}
{\Huge\bfseries\color{uncpblue} SGSI-UNCP\par}
\vspace{0.7cm}
{\LARGE Matriz de Evaluación y Plan de Tratamiento de Riesgos\par}
\vspace{0.5cm}
{\large Documento Base del Sistema de Gestión\\de Seguridad de la Información\par}
\vspace{1.5cm}
{\small Universidad Nacional del Centro del Perú\par}
\vfill
\end{titlepage}
}

\begin{document}

\makecover
\setcounter{tocdepth}{2}
\RaggedRight
\tableofcontents
\newpage

\RaggedRight
\section{Matriz de Evaluación y Plan de Tratamiento de
Riesgos}\label{matriz-de-evaluaciuxf3n-y-plan-de-tratamiento-de-riesgos}

\textbf{Organización:} Universidad Nacional del Centro del Perú (UNCP)
\textbf{Aprobado por:} Oficial de Seguridad y Confianza Digital
\textbf{Metodología:} D-SGSI-04 (NR = O × I) \textbf{Criterio de
Aceptación:} NR \textless= 4 (Bajo)

\begin{center}\rule{0.5\linewidth}{0.5pt}\end{center}

\subsection{Evaluación de Riesgos}\label{evaluaciuxf3n-de-riesgos}

{\setlength{\RaggedRightRightskip}{0pt plus 5em}\small\sloppy \begin{longtable}[]{@{}
  >{\hspace{0pt}\RaggedRight\arraybackslash}p{(\linewidth - 16\tabcolsep) * \real{0.0476}}
  >{\hspace{0pt}\RaggedRight\arraybackslash}p{(\linewidth - 16\tabcolsep) * \real{0.2024}}
  >{\hspace{0pt}\RaggedRight\arraybackslash}p{(\linewidth - 16\tabcolsep) * \real{0.1071}}
  >{\hspace{0pt}\RaggedRight\arraybackslash}p{(\linewidth - 16\tabcolsep) * \real{0.1905}}
  >{\centering\arraybackslash}p{(\linewidth - 16\tabcolsep) * \real{0.0357}}
  >{\centering\arraybackslash}p{(\linewidth - 16\tabcolsep) * \real{0.0357}}
  >{\centering\arraybackslash}p{(\linewidth - 16\tabcolsep) * \real{0.0357}}
  >{\hspace{0pt}\RaggedRight\arraybackslash}p{(\linewidth - 16\tabcolsep) * \real{0.0833}}
  >{\hspace{0pt}\RaggedRight\arraybackslash}p{(\linewidth - 16\tabcolsep) * \real{0.2619}}@{}}
\toprule\noalign{}
\begin{minipage}[b]{\linewidth}\raggedright
ID
\end{minipage} & \begin{minipage}[b]{\linewidth}\raggedright
Activo Afectado
\end{minipage} & \begin{minipage}[b]{\linewidth}\raggedright
Amenaza
\end{minipage} & \begin{minipage}[b]{\linewidth}\raggedright
Vulnerabilidad
\end{minipage} & \begin{minipage}[b]{\linewidth}\centering
O
\end{minipage} & \begin{minipage}[b]{\linewidth}\centering
I
\end{minipage} & \begin{minipage}[b]{\linewidth}\centering
NR
\end{minipage} & \begin{minipage}[b]{\linewidth}\raggedright
Nivel
\end{minipage} & \begin{minipage}[b]{\linewidth}\raggedright
Opción de Tratamiento
\end{minipage} \\
\midrule\noalign{}
\endhead
\bottomrule\noalign{}
\endlastfoot
R01 & ERP ADESA (BD Matrícula, Notas) & Inyección SQL / manipulación de
datos & Código legado no auditado, sin WAF, sin parametrización de
consultas & 2 & 4 & 8 & Medio & Mitigar \\
R02 & Campus Virtual (Moodle 4.1) & Ransomware / cifrado de datos &
Falta de backups inmutables, versión de plugins desactualizada & 2 & 4 &
8 & Medio & Mitigar \\
R03 & Centro Médico (HC digitales) & Acceso no autorizado a datos
sensibles (Ley N.° 29733) & Falta de MFA, autenticación por usuario y
contraseña únicamente & 3 & 4 & 12 & Alto & Mitigar \\
R04 & Red LAN de Facultades & Sniffing / interceptación de tráfico &
Tráfico interno no cifrado, segmentación de red insuficiente & 3 & 3 & 9
& Alto & Mitigar \\
R05 & Storage NAS (Backups) & Falla de hardware / pérdida de respaldos &
Disco único sin RAID, falta de replicación off-site & 2 & 4 & 8 & Medio
& Mitigar \\
R06 & Microsoft 365 (Correo) & Phishing dirigido / suplantación de
identidad & Falta de capacitación en seguridad, ausencia de DMARC/DKIM
avanzado & 3 & 3 & 9 & Alto & Mitigar \\
R07 & Huawei Cloud & Configuración incorrecta de IAM / exposición de
datos & Falta de revisión periódica de accesos cloud, ausencia de Cloud
Security Posture Management & 2 & 4 & 8 & Medio & Mitigar \\
R08 & DSpace (Repositorio) & Pérdida de propiedad intelectual / plagio &
Control de acceso básico sin registro de descargas, ausencia de DRM & 2
& 3 & 6 & Medio & Mitigar \\
R09 & Personal OTI (2 profesionales) & Error humano / mala configuración
& Personal insuficiente para \textasciitilde11,700 usuarios, alta carga
laboral, sin segregación de funciones & 3 & 3 & 9 & Alto & Transferir \\
R10 & Sistema GESDOC & Denegación de servicio (DDoS) & Sin CDN, sin
protecciones anti-DDoS, ancho de banda limitado & 2 & 3 & 6 & Medio &
Mitigar \\
R11 & Servidor ADESA (HW-001) & Corte eléctrico prolongado & UPS con
autonomía limitada (30 min), sin generador eléctrico de respaldo & 2 & 4
& 8 & Medio & Mitigar \\
R12 & Portal Web (WordPress) & Defacement / suplantación institucional &
Plugins sin actualizar, sin WAF, sin monitoreo de integridad de archivos
& 2 & 3 & 6 & Medio & Mitigar \\
R13 & VPN Corporativa & Acceso no autorizado desde equipos personales &
Sin control de dispositivos (BYOD), sin MFA obligatorio & 2 & 4 & 8 &
Medio & Mitigar \\
R14 & SIGA/SIAF (RRHH, Planillas) & Fuga de datos personales de
trabajadores & Acceso privilegiado sin auditoría, datos personales sin
cifrar & 2 & 4 & 8 & Medio & Mitigar \\
R15 & Sede Desconcentrada (Filial) & Robo de equipos / pérdida de
información & Sin inventario actualizado de activos, sin cifrado de
discos & 3 & 3 & 9 & Alto & Mitigar \\
\end{longtable}\normalsize}

\subsection{Plan de Tratamiento de
Riesgos}\label{plan-de-tratamiento-de-riesgos}

{\setlength{\RaggedRightRightskip}{0pt plus 5em}\small\sloppy \begin{longtable}[]{@{}
  >{\hspace{0pt}\RaggedRight\arraybackslash}p{(\linewidth - 12\tabcolsep) * \real{0.0400}}
  >{\hspace{0pt}\RaggedRight\arraybackslash}p{(\linewidth - 12\tabcolsep) * \real{0.2600}}
  >{\hspace{0pt}\RaggedRight\arraybackslash}p{(\linewidth - 12\tabcolsep) * \real{0.2300}}
  >{\hspace{0pt}\RaggedRight\arraybackslash}p{(\linewidth - 12\tabcolsep) * \real{0.1300}}
  >{\hspace{0pt}\RaggedRight\arraybackslash}p{(\linewidth - 12\tabcolsep) * \real{0.1400}}
  >{\hspace{0pt}\RaggedRight\arraybackslash}p{(\linewidth - 12\tabcolsep) * \real{0.0800}}
  >{\hspace{0pt}\RaggedRight\arraybackslash}p{(\linewidth - 12\tabcolsep) * \real{0.1200}}@{}}
\toprule\noalign{}
\begin{minipage}[b]{\linewidth}\raggedright
ID
\end{minipage} & \begin{minipage}[b]{\linewidth}\raggedright
Controles ISO 27001:2022
\end{minipage} & \begin{minipage}[b]{\linewidth}\raggedright
Acción de Tratamiento
\end{minipage} & \begin{minipage}[b]{\linewidth}\raggedright
Responsable
\end{minipage} & \begin{minipage}[b]{\linewidth}\raggedright
Fecha Límite
\end{minipage} & \begin{minipage}[b]{\linewidth}\raggedright
Estado
\end{minipage} & \begin{minipage}[b]{\linewidth}\raggedright
NR Residual
\end{minipage} \\
\midrule\noalign{}
\endhead
\bottomrule\noalign{}
\endlastfoot
R01 & A.8.25 (Desarrollo Seguro), A.8.20 (Seguridad Redes) & Implementar
WAF (ModSecurity), auditar código legacy, migrar a consultas
parametrizadas & OTI (Desarrollo) & Q3-2026 & Pendiente & 4 (Bajo) \\
R02 & A.8.13 (Copias de Seguridad), A.8.8 (Gestión de Vulnerabilidades)
& Activar Cloud Backup inmutable en Huawei Cloud con retención de 30
días. Actualizar plugins de Moodle. & OTI (Sistemas) & Q1-2026 & En
Proceso & 4 (Bajo) \\
R03 & A.5.15 (Control de Acceso), A.8.5 (Autenticación Segura) &
Implementar MFA vía Keycloak / Microsoft Entra ID para acceso al sistema
de salud & OTI (Seguridad) & Q2-2026 & Pendiente & 6 (Medio) \\
R04 & A.8.20 (Seguridad Redes), A.8.22 (Separación de Redes) & Desplegar
segmentación por facultad (VLANs), implementar 802.1X para puertos de
red & OTI (Redes) & Q3-2026 & Pendiente & 4 (Bajo) \\
R05 & A.8.13 (Copias de Seguridad), A.7.10 (Pérdida de Equipos) & Migrar
a RAID 6, implementar replicación a Huawei Cloud, probar restauración
mensual & OTI (Sistemas) & Q2-2026 & Pendiente & 3 (Bajo) \\
R06 & A.6.3 (Concienciación), A.8.7 (Protección contra Malware), A.8.20
(Redes) & Campaña de phishing simulada trimestral, habilitar DMARC/DKIM,
implementar Microsoft Defender for Office 365 & OTI (Seguridad) +
Capacitación & Q2-2026 & Pendiente & 4 (Bajo) \\
R07 & A.5.19 (Proveedores), A.8.29 (Cloud Pública), A.5.22 (Cambios) &
Revisión trimestral de roles IAM, implementar CSPM, auditoría de
configuraciones cloud & OTI (Cloud) & Q2-2026 & En Proceso & 4 (Bajo) \\
R08 & A.5.32 (Propiedad Intelectual), A.8.3 (Control de Acceso), A.8.15
(Registro) & Implementar registro de descargas por usuario, habilitar
DOI, auditar accesos administrativos & OTI (Sistemas) + Investigación &
Q3-2026 & Pendiente & 3 (Bajo) \\
R09 & A.5.3 (Segregación de Tareas), A.6.1 (Antecedentes), A.6.3
(Capacitación) & Contratar 2 profesionales adicionales, externalizar
monitoreo SOC, implementar rotación de funciones & DGA / RRHH & Q4-2026
& Pendiente & 4 (Bajo) \\
R10 & A.8.20 (Redes), A.8.23 (Filtrado Web) & Implementar Cloudflare /
CDN, contratar protección anti-DDoS, aumentar ancho de banda a 500 Mbps
& OTI (Redes) & Q3-2026 & Pendiente & 3 (Bajo) \\
R11 & A.7.11 (Infraestructura), A.7.12 (Cableado) & Adquirir generador
eléctrico de respaldo (50 KVA), extender autonomía UPS a 2 horas & OTI
(Infraestructura) & Q4-2026 & Pendiente & 4 (Bajo) \\
R12 & A.8.8 (Vulnerabilidades), A.8.7 (Malware), A.8.15 (Registro) &
Implementar WAF, monitoreo de integridad de archivos (Tripwire OSSEC),
actualizar WordPress y plugins & OTI (Seguridad) & Q2-2026 & Pendiente &
3 (Bajo) \\
R13 & A.8.5 (Autenticación Segura), A.6.7 (Trabajo Remoto) & Exigir MFA
para toda conexión VPN, implementar control de endpoints (NAC), política
de BYOD & OTI (Seguridad) & Q2-2026 & Pendiente & 3 (Bajo) \\
R14 & A.8.2 (Accesos Privilegiados), A.8.11 (Enmascaramiento), A.8.15
(Registro) & Revisar y auditar accesos privilegiados, cifrar datos
personales en reposo, implementar PAM & OTI (Seguridad) + RRHH & Q3-2026
& Pendiente & 4 (Bajo) \\
R15 & A.7.6 (Equipos fuera de instalaciones), A.8.1 (Dispositivos),
A.7.10 (Pérdida) & Cifrar discos de todas las laptops
(BitLocker/FileVault), actualizar inventario, implementar rastreo GPS &
OTI (Sistemas) & Q2-2026 & Pendiente & 3 (Bajo) \\
\end{longtable}\normalsize}

\subsection{Resumen de Riesgos por
Nivel}\label{resumen-de-riesgos-por-nivel}

{\setlength{\RaggedRightRightskip}{0pt plus 5em}\small\sloppy \begin{longtable}[]{@{}lccl@{}}
\toprule\noalign{}
Nivel & Rango NR & Cantidad & Acción Requerida \\
\midrule\noalign{}
\endhead
\bottomrule\noalign{}
\endlastfoot
Crítico & 13 -- 16 & 0 & Acción inmediata del CGD \\
Alto & 9 -- 12 & 5 & Tratamiento urgente (máx. 90 días) \\
Medio & 5 -- 8 & 10 & Tratamiento planificado (plan anual) \\
Bajo & 1 -- 4 & 0 & Aceptado, monitoreo periódico \\
\textbf{Total} & & \textbf{15} & \\
\end{longtable}\normalsize}

\subsection{Seguimiento y
Actualización}\label{seguimiento-y-actualizaciuxf3n}

\begin{itemize}
\tightlist
\item
  \textbf{Próxima evaluación planificada:} Diciembre 2026
\item
  \textbf{Responsable del seguimiento:} Oficial de Seguridad y Confianza
  Digital
\item
  \textbf{Frecuencia de revisión:} Trimestral (estado de tratamientos) /
  Anual (reevaluación completa)
\item
  \textbf{Informe a:} Comité de Gobierno Digital (CGD) --- Revisión por
  la Dirección (Cláusula 9.3)
\end{itemize}

\begin{center}\rule{0.5\linewidth}{0.5pt}\end{center}

\emph{Nota: El Nivel de Riesgo (NR) = Ocurrencia (O) × Impacto (I). NR
\textgreater= 5 requiere tratamiento documentado. NR \textless= 4 se
considera riesgo aceptable.}

\end{document}
